	\chapter{Enumerazione delle funzioni calcolabili}
	In questo capitolo andremo a trattare l'enumerabilità dei programmi, e quindi dell'insieme delle funzioni calcolabili.  
	
	\section{Enumerazione}
	Un insieme infinito $X$ è enumerabile se esiste una biezione $f: X \to \mathbb{N}$. Se $f$ ed $f^{-1}$ sono anche \textbf{effettivamente calcolabili}, l'insieme $X$ si dirà \textbf{effettivamente enumerabile}.
	
	Nello specifico, \textbf{una enumerazione di un insieme} $X$ è una suriezione $g: \mathbb{N} \to X$, (indicata con $X = \{x_0,x_1,x_2, \dots\}$). Quando $g$ è iniettiva, diremo anche che è un'\textbf{enumerazione senza ripetizioni}. 
	
	\subsection{Alcuni risultati tecnici}
	
	I seguenti insieme sono effettivamente enumerabili:
	
	\begin{enumerate}
		\item $\mathbb{N}\times\mathbb{N}$
		\item $\mathbb{N}^+\times\mathbb{N}^+\times\mathbb{N}^+$
		\item $\bigcup_{k>0}\mathbb{N}^k$
	\end{enumerate}
	
	\paragraph{Dimostrazioni}
	\begin{enumerate}
		\item Consideriamo la biezione $\pi : \mathbb{N}\times\mathbb{N} \to \mathbb{N}$, definita da 
		$$
		\pi (m,n) =_{def} 2^m (2n + 1) -1
		$$
		e cui funzione inversa è data da
		
		$$
		\pi^{-1}(x) = (\pi_1(x), \pi_2(x)) \qquad \text{con} \quad \pi_1(x) =_{def}(x+1), \; \pi_2(x) =_{def}\dfrac{1}{2}\left(\dfrac{x+1}{2^{\pi_1(x)}} - 1\right)
		$$
		
		\item Consideriamo la biezione 
		$$
		\xi : \mathbb{N}^+\times\mathbb{N}^+\times\mathbb{N}^+ \to \mathbb{N} : 
		\xi(m,n,q) =_{def} \pi(\pi(m-1,n-1), q-1)
		$$
		$$
		\xi^{-1}(m,n,q) =_{def} (\pi_1(\pi_1(x)) + 1, \pi_2(\pi_1(x)) + 1, \pi_2(x) + 1)
		$$
		le cui funzioni sono \textbf{effettivamente calcolabili}.
		
		\item Si consideri la biezione $\tau : \bigcup_{k > 0} \mathbb{N}^k \to \mathbb{N}$.
		
		$$
		\tau(a_1, \dots, a_k) =_{def} 2^{a_1} + 2^{a_1 + a_2 + 1} + \dots + 2^{a_1 + a_2 + \dots + a_k + k - 1}-1
		$$
		
		questa è evidentemente una funzione calcolabile, ma per verificare che questa sia una biezione, andremo a sfruttare il fatto che ogni numero naturale ha una e una sola espressione binaria decimale. Noto $x$, possiamo trovare i numeri univoci $k \geq 1$ e i termini  $0 \leq b_1  < b_2 < \dots < b_k$ tale che
		
		$$
		x+1 = 2^{b_1} + 2^{b_2}+ \dots + 2^{b_k}
		$$
		
		dove $b_i = a_1 + \dots + a_i + i -1 \Rightarrow a_{i + 1} = b_{i+1} - b_i -1 \ \forall i: 1\leq i < k$ (e $a_1 = b_1$). La funzione $\tau^{-1}(x)$ è quindi effettivamente calcolabile, confermando la biettività effettiva.
		
	\end{enumerate}
	
	\section{Enumerabilità delle funzioni calcolabili}
	
	Dimostriamo passo passo l'enumerabilità dei programmi, partendo dalla consapevolezza che, un programma qualsiasi, è una sequenza finita di istruzioni. Dimostriamo che l'insieme delle istruzioni è enumerabile.
	
	\paragraph{Notazione}
	\begin{itemize}
		\item $\mathcal I$ : insieme di tutte le istruzioni URM.
		\item $\mathcal P$ : insieme di tutti i programmi URM.
	\end{itemize}
	
	\subsection{Enumerabilità delle istruzioni}
	Dimostriamo l'enumerabilità dell'insieme delle istruzioni $\mathcal I$, definendo una biezione $\beta: \mathcal I \to \mathbb{N}$.
	
	\begin{itemize}
		\item $\beta (Z(n)) = 4(n-1)$
		\item $\beta (S(n)) = 4(n-1) + 1$
		\item $\beta (T(m, n)) = 4\pi(m-1, n-1) + 2$
		\item $\beta (J(m, n, q)) = 4\xi(m, n, q) + 3$
	\end{itemize}
	
	la funzione inversa $\beta^{-1}(x)$, consiste nel trovare i valori $u, r : \; x = 4u + r$, con $r \in \{0,1,2,3\}$. Otteniamo quindi
	
	$$
	\beta^{-1}(x) = \begin{cases}
	Z(u+1) & \text{se } r = 0\\
	S(u+1) & \text{se } r = 1\\
	T(\pi_1(x) + 1, \pi_2(u) + 1) & \text{se } r = 2\\
	J( \xi^{-1}(u)) & \text{se } r = 3\\
	\end{cases}
	$$
	
	ricordiamo che $\xi^{-1}(u) = (m,n,q)$. Esistono quindi delle funzioni calcolabili per una biezione tra $\mathcal I, \mathbb N$, rendendo $\mathcal I$ enumerabile.
	
	\subsection{Enumerabilità dei programmi URM}
	Definiamo adesso una biezione $\gamma : \mathcal{P} \to \mathbb{N}$. Sia $P \in \mathcal P$. Se il programma $P = I_1,I_2,\dots,I_s$, poniamo $\gamma(P) = \tau (\beta(I_1), \beta(I_2), \dots, \beta(I_s))$. È facile verificare che $\gamma, \gamma^{-1}$ sono effettivamente calcolabili, e che quindi $\mathcal P$ è effettivamente enumerabile.
	
	\subsection{Numero di Gödel}
	
	Dato un programma $P \in \mathcal P$, $\gamma(P)$ è detto \textbf{numero di   Gödel} di $P$ (o codice di $P$). Poniamo inoltre $P_n = \gamma^{-1}(n)$.
	
	
	\paragraph{Esempio di calcolo di Numero di Gödel}
	
	$P : T(1,3), S(4), Z(6)$. Si ha:
	\begin{itemize}
		\item $\beta(T(1,3)) = 4\pi(0,2) + 2 = 4(2^0(2\cdot 2 + 1) -1) + 2 = 18$
		\item $\beta(S(4)) = 4\cdot3+1 = 13$
		\item $\beta(Z(6)) = 4\cdot5 = 20$
	\end{itemize}
	$$
	\begin{aligned}
	\gamma(P) &= 2^{18} + 2^{18+13+1} + 2^{18+13+20+2} - 1\\
	&= 2^{18} + 2^{32} + 2^{53} - 1\\
	&= 262144 + 4294967296 + 9007199254740992 - 1\\
	&= 9007203549970431
\end{aligned}
	$$
	
	
	
	
	\section{Enumerabilità delle funzioni calcolabili}
	
	Forniamo delle notazioni utili
	\begin{itemize}
		\item $\phi_a^{(n)}$: la funzione $n$-aria calcolata da $P_a$, ossia $f_{P_n}^{(n)}$.
		\item $W_a^{(n)}$: il dominio di $\phi_a^{(n)}$.
		\item $E_a^{(n)}$: l'immagine di  $\phi_a^{(n)}$.
	\end{itemize}
	
	Se $f$ è una funzione \textbf{unaria calcolabile}, $f = \phi_a$ per qualche $a \in \mathbb{N}$, con $a = \gamma(P)$. Diremo che $a$ è un indice per $f$. Siccome esistono infiniti programmi che calcolano la stessa funzione, la sequenza $\phi_0, \phi_1, \phi_2, \dots$ è una sequenza \textbf{con ripetizioni} che contiene tutte le funzioni unarie calcolabili. Generalizzando anche le arietà, la sequenza $\phi_0^{(n)}, \phi_1^{(n)}, \phi_2^{(n)}, \dots$ è una sequenza \textbf{con ripetizioni} che contiene tutte le funzioni $n$-arie calcolabili.
	
	\subsection{Enumerabilità funzioni calcolabili \textit{n}-arie}
	$\mathcal C_n = \{\phi_a^{(n)}: a \in \mathbb{N}\}$, ossia l'insieme delle funzioni $n$-arie calcolabili, è enumerabile.
	%
	A partire dalla sequenza di funzioni calcolabili $n$-arie $\phi_0^{(n)}, \phi_1^{(n)}, \phi_2^{(n)}, ...$ con ripetizioni, possiamo creare una sequenza senza ripetizioni definita secondo
	
	$$
	\begin{cases}
		f(0) = 0, \\
		f(m+1) = \mu z (\phi_z^{(n)} \neq \phi_{f(0)}^{(n)},  \dots, \phi_{f(m)}^{(n)}).
	\end{cases}
	$$
	%
	Esiste quindi una biezione con $\mathbb{N}$.
	
	
	\subsection{Enumerabilità di tutte le funzioni calcolabili}
	$\mathcal C = \bigcup_{n \geq 1} \mathcal C_n$. Essendo ciascun $\mathcal C_n$ enumerabile, allora $\mathcal{C}$ è enumerabile
	
	
	\section{Diagonalizzazione di Cantor}
	Il metodo di Diagonalizzazione di Cantor permette di costruire funzioni $f$ non calcolabili. Questo stesso metodo è stato utilizzato per dimostrare la non enumerabilità dell'insieme $\mathbb R$.
	
	Data un'enumerazione $\chi_0, \chi_1, \chi_2 \dots$ di funzioni parziali da $\mathbb{N}$ in $\mathbb{N}$, il \textbf{metodo diagonale di Cantor} consente di costruire una funzione $\chi^* : \mathbb{N} \to \mathbb{N}$, diversa da tutte le funzioni $\chi_i$ dell'enumerazione 
	
	$$
	\begin{array}{c|ccccc}
		& 0 & 1 & 2 & 3 & \cdots \\ \hline
		\chi_0 & \chi_0(0) & \chi_0(1) & \chi_0(2) & \chi_0(3) & \cdots \\
		\chi_1 & \chi_1(0) & \chi_1(1) & \chi_1(2) & \chi_1(3) & \cdots \\
		\chi_2 & \chi_2(0) & \chi_2(1) & \chi_2(2) & \chi_2(3) & \cdots \\
		\vdots & \vdots & \vdots & \vdots & \vdots & \ddots
	\end{array}
	$$
	
	essendo però tutte le funzioni calcolabili enumerabili, la funzione $\chi^*$ risulterà non calcolabile. Basterà definire $\chi^*(i) \neq \chi_i(i) \ \forall i \in \mathbb N$.
	
	\subsection{Teorema sull'esistenza di una funzione totale, unaria e non calcolabile}
	
	Esiste una funzione totale, unitaria e non calcolabile.
	
	$$
	f(n)=
	\begin{cases}
	\phi_n(n) +1 &\text{se } \phi_n(n)\downarrow\\
	0 &\text{se } \phi_n(n)\uparrow
	\end{cases}
	$$
	
	\newpage
	
	\section{Teorema \textit{s-m-n }o di parametrizzazione}
	Introduciamo il teorema di parametrizzazione dal generale al particolare.
	
	\paragraph{Concetto: specializzazione di funzione}
	Il concetto è prendere una funzione $f(\vec{x}, \vec{y})$ $n+m$-aria, e crearne una specializzazione $f(\vec{y})$ $m$-aria su dei valori fissati di $\vec{x}$.
	
	\paragraph{Spiegazione intuitiva}
	Sia $f(x,y)$ una funzione calcolabile. Dato $a \in \mathbb{N}$, poniamo 
	
	$$
	f_a(y) =_{def} f(a,y) \qquad \left[ f(x,a) = \phi_{e_x} (y)\right]
	$$
	%
	$f_a$ è calcolabile, quindi $f_a = \phi_{e_a}$ per qualche $e_a$. La funzione $a \to e_a$ è totale, ma questo non significa sia sempre calcolabile. Il teorema di parametrizzazione asserisce che esiste una funzione \textbf{calcolabile} e \textbf{totale} $k$, che mappa $x \to k(x) = e_x$, ossia:
	
	$$
	k:  f_x = \phi_{k(x)}.
	$$
	
	
	\subsection{Teorema \textit{s-m-n } in forma semplice}
	Sia $f(x,y)$ una funzione calcolabile. Allora esiste una funzione totale calcolabile $k(x)$ tale che $f(x,y) = \phi_{k(x)}(y)$.
	
	\paragraph{Dimostrazione} Sia $F$ il programma URM che calcola la funzione $f$. Basta esibire un programma URM $Q_a$ che calcoli $f_a$ per ogni valore di $a$. 
	
	\begin{algorithm}
		\caption{$Q_a$}
		\begin{algorithmic}
			\State $T(1,2)$
			\State $Z(1)$
			\State $S(1)$
			\State $\quad \vdots$ \text{ per $a$ volte.}
			\State $S(1)$
			\State $F$
		\end{algorithmic}
	\end{algorithm}
	
	Ora, $f(a,y)$ è calcolata da $Q_a$. Poniamo $k(x) = \text{codice del programma } Q_x$, ossia il numero di Godel $\gamma(Q_x)$. Esistendo una procedura univoca per calcolare il codice dato un programma, $k(x)$ è totale ($\mathcal{PR}$) ed effettivamente calcolabile, quindi calcolabile, e per costruzione, vale che
	
	$$
	f(x,y) = \phi_{k(x)}(y).
	$$
	
	\subsection{Teorema \textit{s-m-n }}
	Nella versione più generale, il teorema $s-m-n$ dice che esiste una procedura effettiva che permette di calcolare il codice (numero di Godel) di una versione specializzata (con argomenti fissati) di un programma con codice noto. Più formalmente, diremo che: $\forall m, n \geq 1$ esiste una funzione totale calcolabile $(m+1)$-aria $s^m_n(e, \vec{x})$ tale che
	
	$$
	\phi_{s(e,x_1,\dots,x_m)}(y_1,\dots,y_n)
	=
	\phi_e(x_1,\dots,x_m,y_1,\dots,y_n).
	$$

$e$ è il codice del programma ``vecchio'', $\vec{x} = x_1, \dots, x_m$ è l'input su cui specializzarlo, e la funzione $s$ calcola il nuovo numero di Godel. La funzione a destra è una funzione $(m+n)$-aria, mentre quella a sinistra è una funzione $n$-aria
	